% Scoring and leaderboard metrics for the Quant Guild Trading Bot competition.
%
% Build:  pdflatex scoring.tex   (twice, for the cross-references)
%
% This is the normative description of every number the leaderboard shows. It is
% written against the code, not against the problem statement: where the two
% could be read differently, what is here is what `src/auction/scoring.py`,
% `src/auction/engine.py` and `harness/evaluate.py` actually do. Each definition
% names the symbol it comes from so the two can be checked against each other.

% TWO BUILDS, ONE SOURCE.
%
% Variations 3 and 4 are embargoed until mock auction 1, but the scoring is the
% same for all four and participants need it from day one. So this file compiles
% either way, and `scripts/build_docs.sh` produces both:
%
%   scoring-public.pdf   variations 1 and 2 — served now
%   scoring.pdf          all four — served once the admin releases them
%
% `\ifreleased` guards every passage that would give away a V3 or V4 rule. If you
% add one, guard it, and check with:  pdftotext scoring-public.pdf - | grep -i ...
% (tests/test_embargo.py does exactly that.)
\documentclass[11pt,a4paper]{article}

\newif\ifreleased
\ifdefined\RELEASED \releasedtrue \fi

\usepackage[T1]{fontenc}
\usepackage[utf8]{inputenc}
\usepackage[margin=2.4cm]{geometry}
\usepackage{amsmath}
\usepackage{amssymb}
\usepackage{booktabs}
\usepackage{enumitem}
\usepackage{microtype}
\usepackage[dvipsnames]{xcolor}
\usepackage[colorlinks=true,linkcolor=NavyBlue,urlcolor=NavyBlue,
            citecolor=NavyBlue]{hyperref}

\newcommand{\code}[1]{\texttt{\small #1}}
% Bookmarks are plain PDF strings: `\code` and math are not valid there.
\pdfstringdefDisableCommands{\def\code#1{#1}}
\newcommand{\clipf}{\operatorname{clip}}
\newcommand{\sd}{\operatorname{sd}}

\title{\textbf{Scoring and leaderboard metrics}\\[2mm]
       \large Quant Guild Trading Bot competition, IIT Madras}
\author{Odd semester 2026}
\date{}

\begin{document}
\maketitle

\begin{abstract}
\noindent
The board does not rank on profit. It ranks on a \emph{doubly normalised} score:
each block's profit is first divided by that block's hidden maximum value, and
then standardised against the other nineteen bots that played the same block.
This document defines that score and every figure reported beside it, gives the
formula for each, and says which of them are ranking quantities and which are
diagnostics. Section~\ref{sec:mock} covers the mock auctions, which run the same
machinery but do not feed the final standing.
\end{abstract}

\tableofcontents

% =====================================================================
\section{Notation}
\label{sec:notation}

One \emph{game} is a group of $N = 20$ bots playing $T = 2000$ rounds, split
into $B = 4$ blocks of $L = 500$. One \emph{iteration} is a fresh draw of groups
across the whole field; a \emph{showdown} is a set of iterations played back to
back and published as one leaderboard.

\begin{center}
\begin{tabular}{@{}lll@{}}
\toprule
Symbol & Meaning & Where it lives \\
\midrule
$i$ & a bot (index within its group) & \code{key} \\
$b \in \{1,\dots,B\}$ & a block of $L$ rounds & \code{BlockSummary.block} \\
$t$ & a round, $1 \le t \le T$ & \code{obs["round"]} \\
$[m_b, M_b]$ & the block's hidden value bounds & \code{block\_bounds} (secret) \\
$x_{i,t}$ & bot $i$'s private value in round $t$ & \code{obs["x"]} \\
$n_t$ & players still solvent in round $t$ & \code{obs["num\_players"]} \\
$X_t$ & $\max$ value over the \emph{active} players & \code{max\_value} \\
$b_{(1)} \ge b_{(2)} \ge \dots$ & the round's bids, sorted & \code{sorted\_bids} \\
$C_{i,t}$ & bot $i$'s capital after round $t$ & \code{Player.capital} \\
$C^{\text{start}}_{i,b},\, C^{\text{end}}_{i,b}$ & capital at the block's ends &
  \code{BlockSummary} \\
\bottomrule
\end{tabular}
\end{center}

\noindent
$X_t$ is the maximum over the players \emph{active that round}. A bankrupt bot
draws no value, so it contributes nothing to $X_t$ and is not counted in $n_t$.
That is why $n_t$ is handed to every bot: as the field thins, $\mathbb{E}[X_t]$
falls.

% =====================================================================
\section{What happens inside a block}

\subsection{Values}
Every active bot independently draws
\begin{equation}
  x_{i,t} \sim \mathcal{U}[m_b,\, M_b],
\end{equation}
with $[m_b, M_b]$ fixed within a block and redrawn at every boundary. Neither
bound, nor the block index, nor the fact that a boundary has just happened, is
ever reported to a bot. Detecting the regime change is part of the problem.

The bounds themselves are drawn, not chosen. Each block independently takes
\begin{equation}
  m_b \sim \mathcal{U}\{10, 20, \dots, 1000\},
  \qquad
  r_b \sim \mathcal{U}\{100, 200, \dots, 10000\},
  \qquad
  M_b = m_b + r_b,
  \label{eq:bounds}
\end{equation}
uniformly over the two grids --- 100 values each, so $10^4$ possible blocks.
Both grids span two orders of magnitude, and two properties of them do real work:

\begin{itemize}[nosep,leftmargin=1.4em]
\item \textbf{$m_b$ is never zero.} A bot cannot treat $x$ as a draw from
  $\mathcal{U}[0, M]$ and read $M$ off the maximum it has seen. The block
  $[1000, 1100]$, where every value sits within 10\% of every other, is exactly
  as likely as $[10, 10010]$.
\item \textbf{The width is independent of the floor.} Scale and spread move
  separately, so ``how far above the floor am I'' and ``how close to the top am
  I'' are different questions with different answers block to block.
\end{itemize}

\noindent
Each \emph{iteration} draws its own schedule of $B$ blocks from the master seed,
so no two iterations share one; every group inside an iteration plays the same
schedule, which is what makes their standardised scores comparable and what a
balanced or finals iteration seeds on. A hand-written schedule remains possible
(\code{bounds\_mode = "fixed"}) for reproducing a specific run.

\subsection{Capital}
\label{sec:capital}
Capital is redrawn independently for every bot at the start of every block:
\begin{equation}
  \kappa \sim \mathcal{U}[0.5,\, 2.5],
  \qquad
  \boxed{\;C^{\text{start}}_{i,b} = m_b + r_b\,\kappa\;}
  \label{eq:capital}
\end{equation}
with $m_b$ and $r_b$ that block's floor and width from~\eqref{eq:bounds}. What a
bot finished the previous block with does \emph{not} carry over.

Capital is anchored to the floor and scaled by the \emph{width}, not by $M_b$.
That is what keeps $\kappa$ meaning one thing --- ``how many block-widths of
headroom I start with'' --- across blocks that look nothing like each other.
Scaling by $M_b$ instead would hand every bot roughly the same bankroll relative
to the spread of values in a narrow block like $[1000, 1100]$, which is the one
quantity the capital draw exists to vary.

It also needs no patching. Since $m_b \ge 10$ and $r_b \ge 100$, the smallest
possible capital is $10 + 100 \times 0.5 = 60$: comfortably positive, with no
floor term and no additive jitter to keep a small block off zero.

Two consequences shape the whole score. First, within a block some bots start on
half a block-width of headroom and some on two and a half, so a strategy that
behaves identically on a thin bankroll and a fat one is being measured on
exactly that. Second, because capital resets, the four blocks are four
\emph{independent} tests: a bot cannot bank a good block and coast.

\subsection{Legality and bankruptcy}
The legal bid ceiling is the bot's own capital. A bid above it, negative,
\code{NaN}, infinite, not a number, or not returned inside the per-round time
limit is filed as a bid of $0$ for that round.

Capital updates as $C_{i,t} = C_{i,t-1} + \text{payoff}$. A bot with
$C_{i,t} \le 0$ is eliminated \emph{for the remainder of that block} and returns
at the next boundary on a fresh draw from~\eqref{eq:capital}. Bankruptcy is
therefore never fatal, but it forfeits the rest of the block, which is expensive
in exactly the way Section~\ref{sec:survival} describes.

\subsection{Payoffs}
With $X = X_t$ and $b_1 \ge b_2 \ge \dots$ the sorted bids of the round:

\begin{description}[leftmargin=2.2em,style=nextline]
\item[Variation 1 — private value, first price]
  The winner takes $x_i - b_1$, using \emph{the winner's own} value; everyone
  else zero.
\item[Variation 2 — common value, first price]
  The winner takes $X - b_1$; everyone else zero.
\ifreleased
\item[Variation 3 — runner-up penalty]
  \[
    \pi^{(1)} = X - b_1, \qquad
    \pi^{(2)} = -\tfrac12 \max\!\left(X - b_1,\; 0\right),
  \]
  and zero for everyone else. When the winner overpays the runner-up pays
  nothing: the penalty never becomes a reward for the winner's loss.
\item[Variation 4 — funded second price]
  If $b_1 \le X$,
  \[
    \pi^{(1)} = X - b_2, \qquad \pi^{(2)} = X - b_1, \qquad
    \pi^{(k)} = -\frac{w_k}{\sum_{j \in F} w_j}\left(\pi^{(1)} + \pi^{(2)}\right)
    \ \ \text{for } k \in F,
  \]
  with $w_3 = 0.5$, $w_4 = 0.3$, $w_5 = 0.2$ and $F = \{3,4,5\} \cap \{1,\dots,n_t\}$.
  The renormalisation is what keeps the round \emph{exactly zero-sum} when fewer
  than five players are active; with $n_t \le 2$ no penalty is collected. If
  instead $b_1 > X$, rank~1 alone takes $X - b_1$ and nobody else is touched —
  that branch is not zero-sum, by design.
\fi
\end{description}

\ifreleased\else
\noindent
Variations 3 and 4 are released after mock auction 1 and are not described here.
Nothing below depends on which variation is being played: the scoring is
identical for all four.
\fi

\noindent
\textbf{Ties.} Variations 1 and 2 let \emph{every} bidder within $10^{-9}$ of
the top bid win, each collecting the full payoff.\ifreleased{} Variations 3 and 4
need distinct ranks, so ties there are broken uniformly at random.\fi

% =====================================================================
\section{The ranking quantity}
\label{sec:score}

\subsection{Step 1 --- normalised block profit \texorpdfstring{$\pi$}{pi}}
\label{sec:pi}
Raw profit is not comparable across blocks: a block with $M_b = 500$ pays out
roughly fifty times one with $M_b = 10$, and inside a single block bots start
from different capitals. So profit is measured in units of the block's own
hidden maximum:
\begin{equation}
  \boxed{\;
  \pi_{i,b} \;=\; \frac{C^{\text{end}}_{i,b} - C^{\text{start}}_{i,b}}{M_b}
  \;}
  \label{eq:pi}
\end{equation}
with $\pi_{i,b} := 0$ in the degenerate case $M_b \le 0$. A bot that went
bankrupt ends the block on zero capital, so its
$\pi_{i,b} = -C^{\text{start}}_{i,b} / M_b$, which
by~\eqref{eq:capital} is at worst about $-2.5$.

\emph{$\pi$ is the only quantity in the whole system with an absolute meaning.}
Everything downstream is relative. $\pi_{i,b} > 0$ means the strategy made money
that block; $\pi_{i,b} < 0$ means it lost money, however well it placed.

\subsection{Step 2 --- standardisation within the group}
\label{sec:points}
Let $\mathcal{G}$ be the group that played block $b$ together, $|\mathcal{G}| = N$.
Write
\begin{equation}
  \mu_b = \frac{1}{N}\sum_{j \in \mathcal{G}} \pi_{j,b},
  \qquad
  \sigma_b = \sqrt{\frac{1}{N}\sum_{j \in \mathcal{G}} (\pi_{j,b} - \mu_b)^2}
\end{equation}
(the \emph{population} standard deviation --- \code{statistics.pstdev} --- not
the sample one; the twenty bots are the whole population of that block, not a
sample from it). Then
\begin{equation}
  z_{i,b} = \clipf\!\left(\frac{\pi_{i,b} - \mu_b}{\sigma_b},\; -3,\; +3\right),
  \qquad
  \boxed{\; P_{i,b} = 50 + 15\, z_{i,b} \;}
  \label{eq:points}
\end{equation}
so $P_{i,b} \in [5, 95]$, centred on $50$.

Two edge cases are defined explicitly. If $\sigma_b = 0$ --- every bot in the
group scored identically --- then $z_{i,b} := 0$ and everyone gets $50$; nobody
is separated by a block that separated nobody. If the group has a single member,
that member gets $50$.

The $\pm 3$ clip is not cosmetic. Without it a single block in which one bot
happened to draw a large $\kappa$ and win a run of rounds could contribute
several hundred points and decide the competition on variance rather than skill.
Clipping caps any one block's contribution at $95$ and floors it at $5$.

\subsection{Step 3 --- iteration score and total score}
\begin{equation}
  S^{\text{iter}}_{i} = \sum_{b=1}^{B} P_{i,b},
  \qquad
  \boxed{\; \text{\code{score}}_i = \sum_{\text{iterations } r} S^{\text{iter}}_{i,r} \;}
  \label{eq:score}
\end{equation}

\noindent
\textbf{\code{score} is the ranking quantity, and the only one.} With $B = 4$
blocks per iteration and points centred on $50$, an average bot scores about
$200$ per iteration; over a three-iteration showdown, about $600$.

Because \code{score} is a \emph{sum}, a bot that played fewer iterations than the
rest scores lower purely for that. This is what \code{mean\_block\_points}
exists to correct, and it is the documented tie-break
(Section~\ref{sec:ranking}).

% =====================================================================
\section{Every column on the board}
\label{sec:metrics}

Below, $\mathcal{B}_i$ is the multiset of all blocks bot $i$ played across the
whole showdown ($|\mathcal{B}_i| = $ \code{blocks}) and $\mathcal{G}_i$ the set
of games ($|\mathcal{G}_i| = $ \code{games}). All of these are computed in
\code{harness/evaluate.py:aggregate}.

\subsection{\code{score} --- Total score}
Equation~\eqref{eq:score}. Higher is better. This is what the board sorts on.

\subsection{\code{mean\_block\_points} --- Mean block score}
\begin{equation}
  \overline{P}_i = \frac{1}{|\mathcal{B}_i|} \sum_{b \in \mathcal{B}_i} P_{i,b}
  \;\in\; [5, 95]
\end{equation}
\code{score} normalised by how much a bot actually played. It is the tie-break,
and it is the honest number to quote when two bots played different numbers of
games. Read it against $50$: above means the bot beat its group's average block,
below means it did not.

\subsection{\code{mean\_normalised\_profit} --- Mean \texorpdfstring{$\pi$}{pi}}
\label{sec:meanpi}
\begin{equation}
  \boxed{\;
  \overline{\pi}_i = \frac{1}{|\mathcal{B}_i|} \sum_{b \in \mathcal{B}_i} \pi_{i,b}
  \;}
\end{equation}
\textbf{The absolute-profitability check.} Standardisation makes $P$ a purely
relative statement: in a group where every bot loses money, somebody still
scores $95$. $\overline{\pi}$ is the answer to ``but did it actually make
money?''\ --- expressed in units of $M_b$, so $\overline{\pi} = +0.4$ means the
strategy earned, on average, forty percent of the block's maximum value.

A bot can rank first on \code{score} with $\overline{\pi} < 0$. That is not a
contradiction and not a bug: it means the bot lost the least in a field that all
lost. It is, however, exactly the situation a report should notice.\ifreleased{}
In variations 3 and 4 it is common: those are transfer games, so the group's
total $\pi$ cannot be positive in the branch where the winner did not overpay.\fi

\subsection{\code{normalised\_profit\_spread} --- Consistency}
\label{sec:consistency}
\begin{equation}
  \boxed{\;
  s_i = \sqrt{\frac{1}{|\mathcal{B}_i|}
        \sum_{b \in \mathcal{B}_i}\left(\pi_{i,b} - \overline{\pi}_i\right)^2}
  \;}
\end{equation}
The population standard deviation of $\pi$ across \emph{all} blocks the bot
played, in every game. \textbf{Lower is better}, and the board sorts this column
ascending for that reason.

This is the column that separates a strategy from a lottery ticket. Two bots can
share a $\overline{\pi}$ of $+0.3$ with one earning $(0.3, 0.3, 0.3, 0.3)$ and
the other $(1.5, -0.4, 0.3, -0.2)$; they are not the same bot, and the second
one is one bad draw from bankruptcy. Note also that a large $s_i$ is often the
signature of a bot that has \emph{not} detected the block boundary: it is tuned
to one regime and is being carried by the blocks that happen to match it.

The four blocks in an iteration differ deliberately in both scale and width, so
$s_i$ measures robustness to the regime change, not just round-to-round noise.

\subsection{\code{worst\_normalised\_profit} --- Worst block \texorpdfstring{$\pi$}{pi}}
\begin{equation}
  \pi^{\min}_i = \min_{b \in \mathcal{B}_i} \pi_{i,b}
\end{equation}
The tail, reported directly. $s_i$ can be dragged up by an unusually
\emph{good} block, which is not a risk; $\pi^{\min}$ cannot. A bot whose
$\pi^{\min}$ approaches $-C^{\text{start}}_{i,b}/M_b$
(Section~\ref{sec:capital}) lost its entire starting capital in some block,
which is the same statement as \code{survival\_rate} $< 1$.

\subsection{\code{survival\_rate} --- Survival}
\label{sec:survival}
\begin{equation}
  \boxed{\;
  \rho_i = \frac{\#\{b \in \mathcal{B}_i : \text{not bankrupt at the end of } b\}}
                {|\mathcal{B}_i|}
  \;}
\end{equation}
The fraction of blocks the bot finished solvent; $1$ when it never went broke.
A block counts as survived when $C^{\text{end}}_{i,b} > 0$ at the block's last
round --- equivalently, when the bot was never eliminated during it.

Survival is a \emph{diagnostic, not a penalty term}: it appears nowhere
in~\eqref{eq:score}. It does not need to. Going bankrupt forfeits the remainder
of the block, so those rounds earn nothing while the rest of the group keeps
compounding, which drives $\pi_{i,b}$ to its floor and $P_{i,b}$ towards $5$.
The scoring penalises bankruptcy through $\pi$; $\rho_i$ just tells you it
happened, and how often.

Read $\rho_i$ together with $\overline{\pi}_i$. A bot with $\overline{\pi} > 0$
and $\rho < 1$ is one that funds good blocks with catastrophic ones --- the
profile most likely to collapse when the hidden bounds change to ones it has
never seen.

\subsection{Raw profit columns}
Reported for context, and never used to rank. With $\Delta_{i,g}$ the net profit
of bot $i$ in game $g$ (the sum of $C^{\text{end}} - C^{\text{start}}$ over that
game's blocks):
\begin{align}
  \text{\code{mean\_net\_profit}} &= \frac{1}{|\mathcal{G}_i|}\sum_{g} \Delta_{i,g}, &
  \text{\code{std\_net\_profit}}  &= \sqrt{\frac{1}{|\mathcal{G}_i|}\sum_{g}
      \left(\Delta_{i,g} - \overline{\Delta}_i\right)^2}, \\
  \text{\code{best\_net\_profit}} &= \max_g \Delta_{i,g}, &
  \text{\code{worst\_net\_profit}} &= \min_g \Delta_{i,g},
\end{align}
and \code{mean\_final\_capital} is the mean of $C^{\text{end}}$ over the last
block of each game. These are in raw currency units and are therefore dominated
by whichever blocks happened to have the largest $M_b$; that is precisely why
they do not rank anybody.

\subsection{Counters}
\begin{center}
\begin{tabular}{@{}lll@{}}
\toprule
Column & Definition & Note \\
\midrule
\code{games}    & $|\mathcal{G}_i|$ & group-games played this showdown \\
\code{blocks}   & $|\mathcal{B}_i|$ & normally $B \times$ \code{games} \\
\code{wins}     & rounds finishing rank 1 & in V1/V2 every tied top bidder wins \\
\code{errors}   & \code{get\_bid} raised, or returned & each filed as a bid of $0$ \\
                & a non-number, \code{NaN}, $\pm\infty$ or $<0$ & \\
\code{timeouts} & rounds over the time limit & each filed as a bid of $0$ \\
\bottomrule
\end{tabular}
\end{center}

\noindent
An over-capital bid is \emph{not} counted as an error: it is legal to attempt
and is simply filed as $0$. Only malformed returns are counted, which is why a
bot can bleed capital through zero-bids with \code{errors} sitting at $0$.

\noindent
\code{wins} is a diagnostic and a deliberately misleading one if read alone: in
variation 1 a bot that wins every round by overbidding has a large \code{wins}
and a deeply negative $\pi$.\ifreleased{} In variation 4 rank 1 and rank 2 are
\emph{both} paid, so \code{wins} counts only half the times the bot was in the
money.\fi

\subsection{Disqualification}
\code{disqualified} is set by the sandbox, not by the scoring: a policy
violation, a memory kill, or repeated failure to return. A disqualified bot
keeps its computed row --- the numbers stay visible --- but always sorts last,
below every bot that finished legally.

% =====================================================================
\section{Ranking and tie-breaks}
\label{sec:ranking}

Within each variation, rows are ordered by the key
\begin{equation}
  \Big(\;\mathbb{1}[\text{disqualified}],\;\; -\text{\code{score}},\;\;
        -\text{\code{mean\_block\_points}}\;\Big)
\end{equation}
ascending. In words:
\begin{enumerate}[nosep]
  \item every legal bot outranks every disqualified one;
  \item then by total score, highest first;
  \item ties broken by mean block score, highest first --- which is what stops a
        bot that played fewer games from being punished for the scheduler's
        arithmetic rather than its own play.
\end{enumerate}
The four variations are ranked on four separate boards. There is no combined
score across variations; a bot entered in one variation is not comparable to a
bot entered in another, because $\pi$ is standardised against a different field.

% =====================================================================
\section{The tournament}
\label{sec:tournament}

A single 20-bot arrangement is far too noisy to rank on --- who you are grouped
with matters as much as how you play.\ifreleased{} In variations 3 and 4, where
payoffs move between players rather than being created by the auction, it matters
more still.\fi{} So each variation is played over several iterations.

\begin{center}
\begin{tabular}{@{}cll@{}}
\toprule
Iteration & Grouping & What it does \\
\midrule
1 & random & the field shuffled into groups of 20 under seed $S_1$ \\
2 & random, fresh seed & reshuffled independently: new distributions,
                         capitals, opponents \\
3 & strength-balanced & snake-seeded by cumulative points (below) \\
4--5 & finals & the top \code{finals\_size} by cumulative score, head to head \\
\bottomrule
\end{tabular}
\end{center}

\noindent
\textbf{Snake seeding.} Bots are sorted by cumulative score and dealt into $G$
groups in a serpentine: ranks $1 \dots G$ go to groups $1 \dots G$, ranks
$G{+}1 \dots 2G$ go back down $G \dots 1$, and so on. Formally, the bot at
sorted position $p$ (0-based) with $(r, c) = \operatorname{divmod}(p, G)$ lands
in group
\begin{equation}
  g(p) = \begin{cases}
    c & r \text{ even} \\
    G - 1 - c & r \text{ odd.}
  \end{cases}
\end{equation}
This \emph{balances} groups rather than segregating them: every group ends up
with roughly the same average strength, which is what makes standardised scores
comparable across groups.

\textbf{Remainders.} Under random grouping, a trailing group smaller than half
of $N$ is folded into the previous group rather than left to play mostly against
filler bots. Snake seeding needs no such rule: it deals every bot into an
existing group by construction.

\textbf{Filler bots.} A group short of $N$ players is topped up with the three
published sample bots so that the highest-bid race is a real one. Fillers play
for real but never appear on a leaderboard and never contribute a row.

\textbf{Independence.} Each iteration uses its own block schedule
$\{[m_b, M_b]\}$, not merely the same schedule under a new seed, so iteration 2
is a genuinely fresh test rather than a re-roll of the same one. Groups within
an iteration are independent games and are run in parallel.

Final ranking is the sum of the two finals iterations, tie-broken by the
qualification total.

\paragraph{How this maps onto a run.}
\code{grouping} is set \emph{per iteration}, so the table above is one run:
\code{["random", "random", "balanced", "finals", "finals"]}. Iterations play in
order, and each one seeds on the standing after the ones before it --- a
balanced iteration snake-seeds on iterations 1--2's points, and a finals
iteration cuts the field on everything above it. A run whose iterations are all
\code{random} needs no standing at all and does not consult one.

A schedule shorter than \code{iterations} \emph{holds its last entry} rather
than cycling. That is what the default, \code{["random", "balanced"]}, relies
on: iteration 1 draws groups blind because there is nothing to seed on yet, and
every iteration after it is strength-balanced on the standing so far, however
many there are. So the tournament runs in full on \emph{every} showdown ---
practice, mock and final alike --- rather than only when someone remembers to
set it up.

Where a run needs a standing it did not produce itself --- the first balanced
iteration of a fresh board --- it takes the last finished showdown \emph{of its
own kind} (Section~\ref{sec:mock}), never whichever board happened to finish
most recently.

% =====================================================================
\section{Mock auctions}
\label{sec:mock}

Two mock showdowns run before the deadline. They use \emph{the same engine, the
same sandbox, the same scoring, the same aggregation and the same tournament
structure} as the graded run --- there is no separate code path, and a mock is
not a reduced format. It plays the full sequence: iteration 1 on blind groups,
every iteration after it strength-balanced on the standing so far. They differ
only in size, in what is done with the result, and in one recorded fact: every
showdown is stamped with what it was for.

\begin{center}
\begin{tabular}{@{}lll@{}}
\toprule
Kind & What it is & Board \\
\midrule
\code{practice} & the clock, every \code{interval\_minutes} & rewritten by the next tick \\
\code{mock}     & an announced mock auction & archived and published \\
\code{final}    & the run after the deadline & the result \\
\bottomrule
\end{tabular}
\end{center}

\noindent
The stamp is not a label. A \code{balanced} or \code{finals} iteration seeds on
a standing, and where that standing comes from decides who plays whom: seeding
takes the last finished run \emph{of the same kind}, so a mock is seeded by the
previous mock and never by a practice run against throwaway bounds. It is also
what keeps a mock board reachable after the two-hourly clock has replaced the
live one.

\begin{center}
\begin{tabular}{@{}llll@{}}
\toprule
 & \textbf{Mock auction 1} & \textbf{Mock auction 2} & \textbf{Final evaluation} \\
\midrule
Variations   & 1, 2            & \ifreleased 1, 2, 3, 4\else all released\fi
             & \ifreleased 1, 2, 3, 4\else all released\fi \\
Submissions  & optional        & optional          & required \\
Grouping     & \multicolumn{3}{l}{the full tournament: iteration 1 random, the rest seeded} \\
Iterations   & reduced         & reduced           & 3 qualification $+$ 2 finals \\
Bounds       & \multicolumn{3}{l}{drawn from the run's own seed --- never reused} \\
Counts towards the result & no & no                & yes \\
Also does    & \multicolumn{2}{l}{releases variations 3 and 4 (mock 1)} & \\
\bottomrule
\end{tabular}
\end{center}

\noindent
Four things follow, and they are the whole point of running them.

\begin{enumerate}[leftmargin=1.4em]
\item \textbf{A mock score is not a partial final score.} Nothing carries over.
  The graded run starts from an empty board, on block bounds and seeds that have
  never been published, against whatever field submits by the deadline. A mock
  rank is a measurement, not a position.

\item \textbf{The hidden bounds are different, deliberately.} Every iteration of
  every run draws its own blocks from~\eqref{eq:bounds}, off a seed that is not
  published until afterwards. There are $10^4$ possible blocks and no schedule is
  ever reused, so a bot tuned to the $\pi$ it earned in a mock is tuned to a
  distribution that will not appear again --- the same trap the local runner
  warns about, one level up. This is why the bounds are no longer the
  competition's principal secret: the seed is.

\item \textbf{The per-block table is the deliverable, not the rank.} Each
  participant is sent their $(C^{\text{start}}_{b}, C^{\text{end}}_{b}, M_b,
  \pi_b, P_b, \text{bankrupt})$ row for every block. That is where a bot's
  failure mode is legible: a $\pi_b$ that collapses in exactly one block is a
  bot that has not detected the boundary; a low $P_b$ against a healthy $\pi_b$
  is a bot that is profitable but being out-played; a bankruptcy tells you which
  block and which round.

\item \textbf{Mock 1 is also the release gate.} The last two variations are
  published the moment mock auction 1's results are, so their mock field
  necessarily has fewer entrants and less-developed bots than the final one.
  Read a mock 2 rank in a late-released variation with that in mind: the
  standardisation in~\eqref{eq:points} is against \emph{that} field, and a thin
  field of under-tuned bots makes $\sigma_b$ small and $z$ --- and so the whole
  board --- unusually sensitive.
\end{enumerate}

\noindent
Operationally, a mock run is an ordinary showdown triggered from the admin
console and stamped \code{mock}; the practice board that rewrites itself every
two hours is the same thing again, stamped \code{practice}, at a smaller size.
Nothing on either is graded. The only run that is graded is the one after the
final deadline, on settings that are not published until it has finished.

% =====================================================================
\section{Reading the board: three worked cases}

\paragraph{A bot ranked 3rd with \texorpdfstring{$\overline{\pi} = -0.15$}{mean pi = -0.15}.}
It is losing money in absolute terms and still beating most of its group. In V1
or V2 that says the whole field is overbidding and the auction is dissipating its
own surplus.\ifreleased{} In V3 or V4 it is unremarkable: those are transfer
games, and the group's total $\pi$ cannot be positive in the branch where the
winner did not overpay.\fi{} Check $\rho_i$: if it is $1$, the bot is simply
being conservative in a bloodbath, which is a perfectly good strategy.

\paragraph{A bot with \texorpdfstring{$\overline{P} = 62$ and $s = 1.1$}{mean P = 62 and s = 1.1}.}
Strong on average, wildly inconsistent. With four blocks per iteration, an $s$
of $1.1$ against a $\overline{\pi}$ of a few tenths means individual blocks are
swinging by more than the block's own maximum value. Look at $\pi^{\min}$ and
$\rho$: this is usually a bot that is excellent in one regime and bankrupt in
another, and it will not survive a fresh block schedule.

\paragraph{A bot with \texorpdfstring{$\rho = 0.75$}{survival = 0.75} and rank 1.}
It went bankrupt in a quarter of its blocks and still tops the board, so the
blocks it survived were dominant enough to absorb three $P$-values near the
floor. That is a real result, but it is a $3\sigma$-clipped one: the clip
in~\eqref{eq:points} caps how much its good blocks can earn while its bad ones
sit at $5$. It is close to the point where one more bad draw flips the ranking,
and it is the first bot to read by hand.

\end{document}
